\documentclass{amsart}
% For dealing with references we use the comment environment
\usepackage{verbatim}
\newenvironment{reference}{\comment}{\endcomment}
%\newenvironment{reference}{}{}
\newenvironment{slogan}{\comment}{\endcomment}
\newenvironment{history}{\comment}{\endcomment}

% For commutative diagrams we use Xy-pic
\usepackage[all]{xy}

% We use 2cell for 2-commutative diagrams.
\xyoption{2cell}
\UseAllTwocells

% We use multicol for the list of chapters between chapters
\usepackage{multicol}

% This is generally recommended for better output
\usepackage{lmodern}
\usepackage[T1]{fontenc}

% For cross-file-references

% Package for hypertext links:
\usepackage{hyperref}

% For any local file, say "hello.tex" you want to link to please

% Theorem environments.
%
\theoremstyle{plain}
\newtheorem{theorem}[subsection]{Theorem}
\newtheorem{proposition}[subsection]{Proposition}
\newtheorem{lemma}[subsection]{Lemma}

\theoremstyle{definition}
\newtheorem{definition}[subsection]{Definition}
\newtheorem{example}[subsection]{Example}
\newtheorem{exercise}[subsection]{Exercise}
\newtheorem{situation}[subsection]{Situation}

\theoremstyle{remark}
\newtheorem{remark}[subsection]{Remark}
\newtheorem{remarks}[subsection]{Remarks}

\numberwithin{equation}{subsection}

% Macros
%
\def\lim{\mathop{\mathrm{lim}}\nolimits}
\def\colim{\mathop{\mathrm{colim}}\nolimits}
\def\Spec{\mathop{\mathrm{Spec}}}
\def\Hom{\mathop{\mathrm{Hom}}\nolimits}
\def\Ext{\mathop{\mathrm{Ext}}\nolimits}
\def\SheafHom{\mathop{\mathcal{H}\!\mathit{om}}\nolimits}
\def\SheafExt{\mathop{\mathcal{E}\!\mathit{xt}}\nolimits}
\def\Sch{\mathit{Sch}}
\def\Mor{\mathop{\mathrm{Mor}}\nolimits}
\def\Ob{\mathop{\mathrm{Ob}}\nolimits}
\def\Sh{\mathop{\mathit{Sh}}\nolimits}
\def\NL{\mathop{N\!L}\nolimits}
\def\CH{\mathop{\mathrm{CH}}\nolimits}
\def\proetale{{pro\text{-}\acute{e}tale}}
\def\etale{{\acute{e}tale}}
\def\QCoh{\mathit{QCoh}}
\def\Ker{\mathop{\mathrm{Ker}}}
\def\Im{\mathop{\mathrm{Im}}}
\def\Coker{\mathop{\mathrm{Coker}}}
\def\Coim{\mathop{\mathrm{Coim}}}

% Boxtimes
%
\DeclareMathSymbol{\boxtimes}{\mathbin}{AMSa}{"02}

%
% Macros for moduli stacks/spaces
%
\def\QCohstack{\mathcal{QC}\!\mathit{oh}}
\def\Cohstack{\mathcal{C}\!\mathit{oh}}
\def\Spacesstack{\mathcal{S}\!\mathit{paces}}
\def\Quotfunctor{\mathrm{Quot}}
\def\Hilbfunctor{\mathrm{Hilb}}
\def\Curvesstack{\mathcal{C}\!\mathit{urves}}
\def\Polarizedstack{\mathcal{P}\!\mathit{olarized}}
\def\Complexesstack{\mathcal{C}\!\mathit{omplexes}}
% \Pic is the operator that assigns to X its picard group, usage \Pic(X)
% \Picardstack_{X/B} denotes the Picard stack of X over B
% \Picardfunctor_{X/B} denotes the Picard functor of X over B
\def\Pic{\mathop{\mathrm{Pic}}\nolimits}
\def\Picardstack{\mathcal{P}\!\mathit{ic}}
\def\Picardfunctor{\mathrm{Pic}}
\def\Deformationcategory{\mathcal{D}\!\mathit{ef}}

\setlength{\emergencystretch}{3em}
\begin{document}
\title[Stacks Project, Tag 0G7Y]{388 --- Derived Stein splitting near geometrically reduced fibres of proper flat families}
\maketitle
\noindent Copyright (C) 2005--2025 Johan de Jong. Stacks Project authors. Source: \href{https://stacks.math.columbia.edu/tag/0G7Y}{Tag 0G7Y}.

\noindent Editorial difficulty note (not author text). Difficulty H2: The proof must split the zero-degree piece of derived pushforward while ensuring the remaining piece has positive Tor amplitude. It proves fibrewise section surjectivity by passing to an algebraically closed residue field and lifting fibre idempotents, then identifies the zero-degree algebra as finite etale on a neighbourhood..

\noindent Editorial provenance note (not author text). History: excerpt from revision \texttt{a0444\allowbreak 6e57e\allowbreak c1fbc\allowbreak 252a8\allowbreak 71afc\allowbreak ec775\allowbreak 2fb28\allowbreak 07b14}, perfect.tex, lines 8182--8263. Reference presentation was adapted; original-author context and core support blocks were retained or added. The mathematical wording is unchanged. Licensed under GNU FDL 1.2 or later; no invariant sections or cover texts. See ../licenses/COPYING. Context blocks reproduce author definitions and supporting results; they are not additional counted problems.

\begin{lemma}
\label{lemma-proper-flat-geom-red}
Let $f : X \to S$ be a morphism of schemes. Let $s \in S$. Assume
\begin{enumerate}
\item $f$ is proper, flat, and of finite presentation, and
\item the fibre $X_s$ is geometrically reduced.
\end{enumerate}
Then, after replacing $S$ by an open neighbourhood of $s$, there
exists a direct sum decomposition
$Rf_*\mathcal{O}_X = f_*\mathcal{O}_X \oplus P$
in $D(\mathcal{O}_S)$ where $f_*\mathcal{O}_X$ is a finite \'etale
$\mathcal{O}_S$-algebra and
$P$ is a perfect of tor amplitude in $[1, \infty)$.
\end{lemma}

\begin{proof}
The proof of this lemma is similar to the proof of
Lemma \href{https://stacks.math.columbia.edu/tag/0E62}{lemma proper flat h0 (Tag 0E62)}
which we suggest the reader read first.
By cohomology and base change
(Lemma \href{https://stacks.math.columbia.edu/tag/0B91}{lemma flat proper perfect direct image general (Tag 0B91)})
the complex $E = Rf_*\mathcal{O}_X$
is perfect and its formation commutes with arbitrary base change.
This first implies that $E$ has tor amplitude in $[0, \infty)$.

\medskip\noindent
We claim that after replacing $S$ by an open neighbourhood of
$s$ we can find a direct sum decomposition
$E = H^0(E) \oplus \tau_{\geq 1}E$ in $D(\mathcal{O}_S)$
with $\tau_{\geq 1}E$ of tor amplitude in $[1, \infty)$.
Assume the claim is true for now and assume we've made the replacement
so we have the direct sum decomposition.
Since $E$ has tor amplitude in $[0, \infty)$ we find that
$H^0(E)$ is a flat $\mathcal{O}_S$-module.
Hence $H^0(E)$ is a flat, perfect $\mathcal{O}_S$-module,
hence finite locally free, see
More on Algebra, Lemma \href{https://stacks.math.columbia.edu/tag/0658}{lemma perfect (Tag 0658)}
(and the fact that finite projective modules are finite locally free by
Algebra, Lemma \href{https://stacks.math.columbia.edu/tag/00NX}{lemma finite projective (Tag 00NX)}).
Of course $H^0(E) = f_*\mathcal{O}_X$ is an $\mathcal{O}_S$-algebra.
By cohomology and base change we obtain
$H^0(E) \otimes \kappa(s) = H^0(X_s, \mathcal{O}_{X_s})$.
By Varieties, Lemma
\href{https://stacks.math.columbia.edu/tag/0BUG}{lemma proper geometrically reduced global sections (Tag 0BUG)}
and the assumption that $X_s$ is geometrically reduced, we
see that $\kappa(s) \to H^0(E) \otimes \kappa(s)$
is finite \'etale. By Morphisms, Lemma
\href{https://stacks.math.columbia.edu/tag/0476}{lemma set points where fibres etale (Tag 0476)}
applied to the finite locally free morphism
$\underline{\Spec}_S(H^0(E)) \to S$,
we conclude that after shrinking $S$ the $\mathcal{O}_S$-algebra
$H^0(E)$ is finite \'etale.

\medskip\noindent
It remains to prove the claim. For this it suffices to prove that the map
$$
(f_*\mathcal{O}_X)_s
\longrightarrow
H^0(X_s, \mathcal{O}_{X_s}) = H^0(E \otimes^\mathbf{L} \kappa(s))
$$
is surjective, see
More on Algebra, Lemma \href{https://stacks.math.columbia.edu/tag/0A1U}{lemma better cut complex in two (Tag 0A1U)}.
Choose a flat local ring homomorphism $\mathcal{O}_{S, s} \to A$
such that the residue field $k$ of $A$ is algebraically closed, see
Algebra, Lemma \href{https://stacks.math.columbia.edu/tag/03C3}{lemma flat local given residue field (Tag 03C3)}.
By flat base change (Cohomology of Schemes, Lemma
\href{https://stacks.math.columbia.edu/tag/02KH}{lemma flat base change cohomology (Tag 02KH)})
we get $H^0(X_A, \mathcal{O}_{X_A}) =
(f_*\mathcal{O}_X)_s \otimes_{\mathcal{O}_{S, s}} A$
and $H^0(X_k, \mathcal{O}_{X_k}) =
H^0(X_s, \mathcal{O}_{X_s}) \otimes_{\kappa(s)} k$.
Hence it suffices to prove that
$H^0(X_A, \mathcal{O}_{X_A}) \to H^0(X_k, \mathcal{O}_{X_k})$
is surjective. Since $X_k$ is a reduced proper scheme over $k$
and since $k$ is algebraically closed, we see that
$H^0(X_k, \mathcal{O}_{X_k})$ is a finite product of copies
of $k$ by the already used Varieties, Lemma
\href{https://stacks.math.columbia.edu/tag/0BUG}{lemma proper geometrically reduced global sections (Tag 0BUG)}.
Since by Lemma \href{https://stacks.math.columbia.edu/tag/0G7X}{lemma proper idempotent on fibre (Tag 0G7X)}
the idempotents of this $k$-algebra are in the image
of $H^0(X_A, \mathcal{O}_{X_A}) \to H^0(X_k, \mathcal{O}_{X_k})$ we conclude.
\end{proof}
\end{document}
